\chapter[Le tri]{
Le tri}
{
Dans ce chapitre nous voyons quelques algorithmes simples pour trier un
ensemble d'informations : recherche des maxima, tri
par insertion et tri bulle. Des algorithmes plus efficaces seront vus
en deuxième année.}

\begin{center}
 [Warning: Image ignored] % Unhandled or unsupported graphics:
%\includegraphics[width=1.27cm,height=1.27cm]{log1-img/log1-img115}

\end{center}
\section{Motivation}
{
Une application importante de l’informatique est le tri d’informations.
La conception d’algorithmes efficaces pour ce type de traitement s’est
imposée avec l’apparition de bases de données de taille importante. On
citera par exemple la recherche d’un numéro de téléphone ou une adresse
dans la version informatisée des pages blanches, ou de façon plus
spectaculaire l’efficacité de certains moteurs de recherche qui
permettent de retrouver en quelques secondes des mots clés arbitraires
parmi plusieurs millions de pages sur le web.}

{
La recherche efficace d’information implique un tri préalable de
celle-ci. En effet, si l’information n’est pas classée ou triée, le
seul algorithme possible reviendrait à parcourir entièrement l’ensemble
des informations. Pour exemple, il suffit d’imaginer un dictionnaire
dans lequel les mots seraient mélangés de façon aléatoire au lieu
d’être classés par ordre alphabétique. Pour trouver le moindre mot dans
ce dictionnaire, il faudrait à chaque fois le parcourir entièrement !
Il est clair que le classement préalable (ordre alphabétique) accélère
grandement la recherche.}

{
Ainsi, recherche et tri sont étroitement liés, et la façon dont les
informations sont triées conditionne bien entendu la façon de
rechercher l’information (cf. algorithme de recherche dichotomique).
Pour exemple, prenons cette fois-ci un dictionnaire des mots croisés
dans lequel les mots sont d’abord regroupés selon leur longueur et
ensuite par ordre alphabétique. La façon de rechercher un mot dans ce
dictionnaire est bien sûr différente de la recherche dans un
dictionnaire usuel. }

{
Le problème central est donc le tri des informations. Celui-ci a pour
but d’organiser un ensemble d’informations qui ne l’est pas \textit{à
priori. }On peut distinguer trois grands cas de figure :}

\liststyleWWviiiNumi
\begin{enumerate}
\item {
D’abord les situations impliquant un classement total d’un ensemble de
données «~brutes~», c’est-à-dire complètement désordonnées. Prenons
pour exemple les feuilles récoltées en vrac à l’issue d’un examen ; il
y a peu de chances que celles-ci soient remises à l’examinateur de
manière ordonnée ; celui-ci devra donc procéder au tri de l’ensemble
des copies, par exemple par ordre alphabétique des noms des étudiants,
ou par numéro de groupe etc.}
\item {
Ensuite les situations où on s’arrange pour ne jamais devoir trier la
totalité des éléments d’un ensemble, qui resterait cependant à tout
moment ordonné. Imaginons le cas d’une bibliothèque~dont les livres
sont rangés par ordre alphabétique des auteurs : à l’achat d’un nouveau
livre, ou au retour de prêt d’un livre, celui-ci est immédiatement
rangé à la bonne place. Ainsi, l’ordre global de la bibliothèque est
maintenu par la répétition d’une seule opération élémentaire consistant
à insérer à la bonne place un livre parmi la collection. C’est la
situation que nous avons considérée avec la liste ordonnée.}
\item {
Enfin, les situations qui consistent à devoir re-trier des données
préalablement ordonnées sur un autre critère. Prenons l’exemple d’un
paquet de copies d’examen déjà triées sur l’ordre alphabétique des noms
des étudiants, et qu’on veut re-trier cette fois-ci sur les numéros de
groupe. Il est clair qu’une méthode efficace veillera à conserver
l’ordre alphabétique déjà présent dans la première situation afin que
les copies apparaissent dans cet ordre dans chacun des groupes.}
\end{enumerate}
{
Le dernier exemple illustre un classement sur une \textbf{clé complexe}
(ou \textbf{composée}) impliquant la comparaison de plusieurs champs
d’une même structure : le premier classement se fait sur le numéro de
groupe, et à numéro de groupe égal, l’ordre se départage sur le nom de
l’étudiant. On dira de cet ensemble qu’il est classé en \textbf{majeur}
sur le numéro de groupe et en \textbf{mineur} sur le nom d’étudiant.}

{
Notons que certains tris sont dits \textbf{stables} parce
qu'en cas de tri sur une nouvelle clé, l’ordre de la
clé précédente est préservé pour des valeurs identiques de la nouvelle
clé, ce qui évite de faire des comparaisons sur les deux champs à la
fois. Les méthodes nommées \textbf{tri par insertion}, \textbf{tri
bulle} et \textbf{tri par recherche de minima successifs }(que nous
allons aborder ci-après) sont stables.}

{
Le tri d’un ensemble d’informations n’offre que des avantages. Outre le
fait déjà mentionné de permettre une recherche et une obtention rapide
de l’information, il permet aussi l’application de traitements
algorithmiques efficaces (comme par exemple celui des ruptures que nous
verrons plus loin) qui s’avéreraient trop coûteux (en temps, donc en
argent !) s’ils s’effectuaient sur des ensembles non triés.}

{
Certains tris sont évidemment plus performants que d’autres. Le choix
d’un tri particulier dépend de la taille de l’ensemble à trier et de la
manière dont il se présente, c’est-à-dire déjà plus ou moins ordonné.
La performance quant à elle se juge sur deux critères : le nombre de
tests effectués (comparaisons de valeurs) et le nombre de transferts de
valeurs réalisés. }

{
Les algorithmes classiques de tri présentés dans ce chapitre le sont
surtout à titre pédagogique. Ils ont tous une «~complexité en
\textit{n}²~», ce qui veut dire que si \textit{n} est le nombre
d’éléments à trier, le nombre d’opérations élémentaires (tests et
transferts de valeurs) est proportionnel à \textit{n}². Ils conviennent
donc pour des ensembles de «~taille raisonnable~», mais peuvent devenir
extrêmement lents à l’exécution pour le tri de grands ensembles, comme
par exemple les données de l’annuaire téléphonique. Plusieurs solutions
existent, comme la méthode de tri \textbf{Quicksort}. Cet algorithme
très efficace faisant appel à la récursivité et qui sera étudié en
deuxième année a une complexité en \textit{n} log \textit{n}. }

{
Dans ce chapitre, les algorithmes traiteront du tri dans un
\textbf{tableau d’entiers à une }\textbf{dimension}. Toute autre
situation peut bien entendu se ramener à celle-ci moyennant la
définition de la relation d’ordre propre au type de données utilisé. Ce
sera par exemple, l’ordre alphabétique pour les chaines de caractères,
l’ordre chronologique pour des objets \textstyleCodeInsr{Date} ou
\textstyleCodeInsr{Moment }\textstyleCodeInsr{\textrm{(que nous verrons
plus tard)}}, etc. De plus, le seul ordre envisagé sera l’ordre
\textbf{croissant} des données. On peut tout aussi bien envisager le
tri d’une liste en lieu et place d’un tableau.}

\begin{center}
 [Warning: Image ignored] % Unhandled or unsupported graphics:
%\includegraphics[width=1.323cm,height=1.323cm]{log1-img/log1-img116}

\end{center}
{
Enfin, dans toutes les méthodes de tri abordées, nous supposerons la
taille physique du tableau à trier (notée \textit{n}) égale à sa taille
logique, celle-ci n’étant pas modifiée par l’action de tri.}

\section{Tri par insertion}
{
Cette méthode de tri repose sur le principe d’insertion de valeurs dans
un tableau ordonné que nous allons développer lors de la description de
la liste ordonnée. }

{\sffamily\bfseries\upshape
Description de l’algorithme}

{
Le tableau à trier sera à chaque étape subdivisé en deux sous-tableaux :
le premier cadré à gauche contiendra des éléments déjà ordonnés, et le
second, cadré à droite, ceux qu’il reste à insérer dans le sous-tableau
trié. Celui-ci verra sa taille s’accroitre au fur et à mesure des
insertions, tandis que celle du sous-tableau des éléments non triés
diminuera progressivement.}

{
Au départ de l’algorithme, le sous-tableau trié est le premier élément
du tableau. Comme il ne possède qu’un seul élément, ce sous-tableau est
donc bien ordonné ! Chaque étape consiste ensuite à prendre le premier
élément du sous-tableau non trié et à l’insérer à la bonne place dans
le sous-tableau trié.}

{
Prenons comme exemple un tableau tab de 20 entiers. Au départ, le
sous-tableau trié est formé du premier élément tab[1] = 20 :}


\bigskip

\begin{center}
\tablehead{}
\begin{supertabular}{m{0.49700004cm}m{0.49700004cm}m{0.49700004cm}m{0.49700004cm}m{0.49700004cm}m{0.49700004cm}m{0.49700004cm}m{0.49700004cm}m{0.49700004cm}m{0.49700004cm}m{0.49700004cm}m{0.49700004cm}m{0.49700004cm}m{0.49700004cm}m{0.49700004cm}m{0.49700004cm}m{0.49700004cm}m{0.48300004cm}m{0.48300004cm}m{0.53000003cm}}
\centering \sffamily\itshape 1 &
\centering \sffamily\itshape 2 &
\centering \sffamily\itshape 3 &
\centering \sffamily\itshape 4 &
\centering \sffamily\itshape 5 &
\centering \sffamily\itshape 6 &
\centering \sffamily\itshape 7 &
\centering \sffamily\itshape 8 &
\centering \sffamily\itshape 9 &
\centering \sffamily\itshape 10 &
\centering \sffamily\itshape 11 &
\centering \sffamily\itshape 12 &
\centering \sffamily\itshape 13 &
\centering \sffamily\itshape 14 &
\centering \sffamily\itshape 15 &
\centering \sffamily\itshape 16 &
\centering \sffamily\itshape 17 &
\centering \sffamily\itshape 18 &
\centering \sffamily\itshape 19 &
\centering\arraybslash \sffamily\itshape
20\\\hline
\multicolumn{1}{|m{0.49700004cm}|}{ 20} &
\multicolumn{1}{m{0.49700004cm}|}{ 12} &
\multicolumn{1}{m{0.49700004cm}|}{ 18} &
\multicolumn{1}{m{0.49700004cm}|}{ 17} &
\multicolumn{1}{m{0.49700004cm}|}{ 15} &
\multicolumn{1}{m{0.49700004cm}|}{ 14} &
\multicolumn{1}{m{0.49700004cm}|}{ 15} &
\multicolumn{1}{m{0.49700004cm}|}{ 16} &
\multicolumn{1}{m{0.49700004cm}|}{ 18} &
\multicolumn{1}{m{0.49700004cm}|}{ 17} &
\multicolumn{1}{m{0.49700004cm}|}{ 12} &
\multicolumn{1}{m{0.49700004cm}|}{ 14} &
\multicolumn{1}{m{0.49700004cm}|}{ 16} &
\multicolumn{1}{m{0.49700004cm}|}{ 18} &
\multicolumn{1}{m{0.49700004cm}|}{ 15} &
\multicolumn{1}{m{0.49700004cm}|}{ 15} &
\multicolumn{1}{m{0.49700004cm}|}{ 19} &
\multicolumn{1}{m{0.48300004cm}|}{ 11} &
\multicolumn{1}{m{0.48300004cm}|}{ 11} &
\multicolumn{1}{m{0.53000003cm}|}{ 13}\\\hline
\end{supertabular}
\end{center}

\bigskip

{
L’étape suivante consiste à insérer tab[2] = 12 dans ce sous-tableau, de
taille 2 :}


\bigskip

\begin{center}
\tablehead{}
\begin{supertabular}{m{0.49700004cm}m{0.49700004cm}m{0.49700004cm}m{0.49700004cm}m{0.49700004cm}m{0.49700004cm}m{0.49700004cm}m{0.49700004cm}m{0.49700004cm}m{0.49700004cm}m{0.49700004cm}m{0.49700004cm}m{0.49700004cm}m{0.49700004cm}m{0.49700004cm}m{0.49700004cm}m{0.49700004cm}m{0.48300004cm}m{0.48300004cm}m{0.53000003cm}}
\centering \sffamily\itshape 1 &
\centering \sffamily\itshape 2 &
\centering \sffamily\itshape 3 &
\centering \sffamily\itshape 4 &
\centering \sffamily\itshape 5 &
\centering \sffamily\itshape 6 &
\centering \sffamily\itshape 7 &
\centering \sffamily\itshape 8 &
\centering \sffamily\itshape 9 &
\centering \sffamily\itshape 10 &
\centering \sffamily\itshape 11 &
\centering \sffamily\itshape 12 &
\centering \sffamily\itshape 13 &
\centering \sffamily\itshape 14 &
\centering \sffamily\itshape 15 &
\centering \sffamily\itshape 16 &
\centering \sffamily\itshape 17 &
\centering \sffamily\itshape 18 &
\centering \sffamily\itshape 19 &
\centering\arraybslash \sffamily\itshape
20\\\hline
\multicolumn{1}{|m{0.49700004cm}|}{ 12} &
\multicolumn{1}{m{0.49700004cm}|}{ 20} &
\multicolumn{1}{m{0.49700004cm}|}{ 18} &
\multicolumn{1}{m{0.49700004cm}|}{ 17} &
\multicolumn{1}{m{0.49700004cm}|}{ 15} &
\multicolumn{1}{m{0.49700004cm}|}{ 14} &
\multicolumn{1}{m{0.49700004cm}|}{ 15} &
\multicolumn{1}{m{0.49700004cm}|}{ 16} &
\multicolumn{1}{m{0.49700004cm}|}{ 18} &
\multicolumn{1}{m{0.49700004cm}|}{ 17} &
\multicolumn{1}{m{0.49700004cm}|}{ 12} &
\multicolumn{1}{m{0.49700004cm}|}{ 14} &
\multicolumn{1}{m{0.49700004cm}|}{ 16} &
\multicolumn{1}{m{0.49700004cm}|}{ 18} &
\multicolumn{1}{m{0.49700004cm}|}{ 15} &
\multicolumn{1}{m{0.49700004cm}|}{ 15} &
\multicolumn{1}{m{0.49700004cm}|}{ 19} &
\multicolumn{1}{m{0.48300004cm}|}{ 11} &
\multicolumn{1}{m{0.48300004cm}|}{ 11} &
\multicolumn{1}{m{0.53000003cm}|}{ 13}\\\hline
\end{supertabular}
\end{center}

\bigskip

{
Ensuite, c’est au tour de tab[3] = 18 d’être inséré :}


\bigskip

\begin{center}
\tablehead{}
\begin{supertabular}{m{0.499cm}m{0.499cm}m{0.499cm}m{0.499cm}m{0.499cm}m{0.499cm}m{0.499cm}m{0.499cm}m{0.499cm}m{0.499cm}m{0.499cm}m{0.499cm}m{0.499cm}m{0.499cm}m{0.499cm}m{0.499cm}m{0.499cm}m{0.48400003cm}m{0.48400003cm}m{0.49700004cm}}
\centering \sffamily\itshape 1 &
\centering \sffamily\itshape 2 &
\centering \sffamily\itshape 3 &
\centering \sffamily\itshape 4 &
\centering \sffamily\itshape 5 &
\centering \sffamily\itshape 6 &
\centering \sffamily\itshape 7 &
\centering \sffamily\itshape 8 &
\centering \sffamily\itshape 9 &
\centering \sffamily\itshape 10 &
\centering \sffamily\itshape 11 &
\centering \sffamily\itshape 12 &
\centering \sffamily\itshape 13 &
\centering \sffamily\itshape 14 &
\centering \sffamily\itshape 15 &
\centering \sffamily\itshape 16 &
\centering \sffamily\itshape 17 &
\centering \sffamily\itshape 18 &
\centering \sffamily\itshape 19 &
\centering\arraybslash \sffamily\itshape
20\\\hline
\multicolumn{1}{|m{0.499cm}|}{ 12} &
\multicolumn{1}{m{0.499cm}|}{ 18} &
\multicolumn{1}{m{0.499cm}|}{ 20} &
\multicolumn{1}{m{0.499cm}|}{ 17} &
\multicolumn{1}{m{0.499cm}|}{ 15} &
\multicolumn{1}{m{0.499cm}|}{ 14} &
\multicolumn{1}{m{0.499cm}|}{ 15} &
\multicolumn{1}{m{0.499cm}|}{ 16} &
\multicolumn{1}{m{0.499cm}|}{ 18} &
\multicolumn{1}{m{0.499cm}|}{ 17} &
\multicolumn{1}{m{0.499cm}|}{ 12} &
\multicolumn{1}{m{0.499cm}|}{ 14} &
\multicolumn{1}{m{0.499cm}|}{ 16} &
\multicolumn{1}{m{0.499cm}|}{ 18} &
\multicolumn{1}{m{0.499cm}|}{ 15} &
\multicolumn{1}{m{0.499cm}|}{ 15} &
\multicolumn{1}{m{0.499cm}|}{ 19} &
\multicolumn{1}{m{0.48400003cm}|}{ 11} &
\multicolumn{1}{m{0.48400003cm}|}{ 11} &
\multicolumn{1}{m{0.49700004cm}|}{ 13}\\\hline
\end{supertabular}
\end{center}

\bigskip

{
Et ainsi de suite jusqu’à insertion du dernier élément dans le
sous-tableau trié. }

{\sffamily\bfseries
Algorithme}

{
On combine la recherche de la position d’insertion et le décalage
concomitant de valeurs.}

{\sffamily
\textstyleMotCl{module} triInsertion ( tab↓↑ : tableau[1 à n] d’entiers
)}


\bigskip

{\sffamily
\ \ i, j, valAInsérer : entiers}

{\sffamily
\ \ \textstyleMotCl{pour} i \textstyleMotCl{de} 2 \textstyleMotCl{à} n
\textstyleMotCl{faire}}

{\sffamily
\ \ \ \ valAInsérer {\textrm{[F0DF?]}} tab[i]}

{\sffamily
\ \ \ \ // Recherche de l’endroit où insérer valAInsérer dans le }

{\sffamily
\ \ \ \ // sous-tableau trié et décalage simultané des éléments}

{\sffamily
\ \ \ \ j \textrm{[F0DF?]} i - 1}

{\sffamily
\ \ \ \ \textstyleMotCl{tant} \textstyleMotCl{que} j ${\geq}$ 1 ET
valAInsérer {\textless} tab[j] \textstyleMotCl{faire}}

{\sffamily
\ \ \ \ \ \ {tab[j+1]
}\textrm{[F0DF?]}{ tab[j]}}

{\sffamily
{\ \ \ \ \ \ j
}\textrm{[F0DF?]}{ j – 1}}

{\sffamily
{\ \ \ \ }\textstyleMotCl{fin}
\textstyleMotCl{tant} \textstyleMotCl{que}}

{\sffamily
\ \ \ \ tab[j+1] \textrm{[F0DF?]} \ valAInsérer}

{\sffamily
\ \ \textstyleMotCl{fin} \textstyleMotCl{pour}}


\bigskip

{\sffamily
\textstyleMotCl{fin} \textstyleMotCl{module}}


\bigskip

\section{Tri par sélection des minima successifs}
{
Dans ce tri, on recherche à chaque étape la plus petite valeur de
l’ensemble non encore trié et on peut la placer à sa position
définitive.}

{\sffamily\bfseries\upshape
Description de l’algorithme}

{
Prenons par exemple un tableau de 20 entiers. La première étape consiste
à rechercher la valeur minimale du tableau. Il s’agit de l’élément
d’indice 10 et de valeur 17.}

\begin{center}
\tablehead{}
\begin{supertabular}{m{0.49700004cm}m{0.49700004cm}m{0.49700004cm}m{0.49700004cm}m{0.49700004cm}m{0.49700004cm}m{0.49700004cm}m{0.49700004cm}m{0.49700004cm}m{0.49700004cm}m{0.49700004cm}m{0.49700004cm}m{0.49700004cm}m{0.49700004cm}m{0.49700004cm}m{0.49700004cm}m{0.49700004cm}m{0.49700004cm}m{0.49700004cm}m{0.532cm}}
\itshape 1 &
\itshape 2 &
\itshape 3 &
\itshape 4 &
\itshape 5 &
\itshape 6 &
\itshape 7 &
\itshape 8 &
\itshape 9 &
\itshape 10 &
\itshape 11 &
\itshape 12 &
\itshape 13 &
\itshape 14 &
\itshape 15 &
\itshape 16 &
\itshape 17 &
\itshape 18 &
\itshape 19 &
\itshape 20\\\hline
\multicolumn{1}{|m{0.49700004cm}|}{ 20} &
\multicolumn{1}{m{0.49700004cm}|}{ 52} &
\multicolumn{1}{m{0.49700004cm}|}{ 61} &
\multicolumn{1}{m{0.49700004cm}|}{ 47} &
\multicolumn{1}{m{0.49700004cm}|}{ 82} &
\multicolumn{1}{m{0.49700004cm}|}{ 64} &
\multicolumn{1}{m{0.49700004cm}|}{ 95} &
\multicolumn{1}{m{0.49700004cm}|}{ 66} &
\multicolumn{1}{m{0.49700004cm}|}{ 84} &
\multicolumn{1}{m{0.49700004cm}|}{ 17} &
\multicolumn{1}{m{0.49700004cm}|}{ 32} &
\multicolumn{1}{m{0.49700004cm}|}{ 24} &
\multicolumn{1}{m{0.49700004cm}|}{ 46} &
\multicolumn{1}{m{0.49700004cm}|}{ 48} &
\multicolumn{1}{m{0.49700004cm}|}{ 75} &
\multicolumn{1}{m{0.49700004cm}|}{ 55} &
\multicolumn{1}{m{0.49700004cm}|}{ 19} &
\multicolumn{1}{m{0.49700004cm}|}{ 61} &
\multicolumn{1}{m{0.49700004cm}|}{ 21} &
\multicolumn{1}{m{0.532cm}|}{ 30}\\\hline
\end{supertabular}
\end{center}
{
Celui-ci devrait donc apparaitre en 1\textsuperscript{ère }position du
tableau. Hors, cette position est occupée par la valeur 20. Comment
procéder dès lors pour ne perdre aucune valeur et sans faire appel à un
second tableau ? La solution est simple, il suffit d’échanger le
contenu des deux éléments d’indices 1 et 10 :}

\begin{center}
\tablehead{}
\begin{supertabular}{m{0.49700004cm}m{0.49700004cm}m{0.49700004cm}m{0.49700004cm}m{0.49700004cm}m{0.49700004cm}m{0.49700004cm}m{0.49700004cm}m{0.49700004cm}m{0.49700004cm}m{0.49700004cm}m{0.49700004cm}m{0.49700004cm}m{0.49700004cm}m{0.49700004cm}m{0.49700004cm}m{0.49700004cm}m{0.49700004cm}m{0.49700004cm}m{0.532cm}}
\itshape 1 &
\itshape 2 &
\itshape 3 &
\itshape 4 &
\itshape 5 &
\itshape 6 &
\itshape 7 &
\itshape 8 &
\itshape 9 &
\itshape 10 &
\itshape 11 &
\itshape 12 &
\itshape 13 &
\itshape 14 &
\itshape 15 &
\itshape 16 &
\itshape 17 &
\itshape 18 &
\itshape 19 &
\itshape 20\\\hline
\multicolumn{1}{|m{0.49700004cm}|}{ 17} &
\multicolumn{1}{m{0.49700004cm}|}{ 52} &
\multicolumn{1}{m{0.49700004cm}|}{ 61} &
\multicolumn{1}{m{0.49700004cm}|}{ 47} &
\multicolumn{1}{m{0.49700004cm}|}{ 82} &
\multicolumn{1}{m{0.49700004cm}|}{ 64} &
\multicolumn{1}{m{0.49700004cm}|}{ 95} &
\multicolumn{1}{m{0.49700004cm}|}{ 66} &
\multicolumn{1}{m{0.49700004cm}|}{ 84} &
\multicolumn{1}{m{0.49700004cm}|}{ 20} &
\multicolumn{1}{m{0.49700004cm}|}{ 32} &
\multicolumn{1}{m{0.49700004cm}|}{ 24} &
\multicolumn{1}{m{0.49700004cm}|}{ 46} &
\multicolumn{1}{m{0.49700004cm}|}{ 48} &
\multicolumn{1}{m{0.49700004cm}|}{ 75} &
\multicolumn{1}{m{0.49700004cm}|}{ 55} &
\multicolumn{1}{m{0.49700004cm}|}{ 19} &
\multicolumn{1}{m{0.49700004cm}|}{ 61} &
\multicolumn{1}{m{0.49700004cm}|}{ 21} &
\multicolumn{1}{m{0.532cm}|}{ 30}\\\hline
\end{supertabular}
\end{center}
{
Le tableau se subdivise à présent en deux sous-tableaux, un sous-tableau
déjà trié (pour le moment réduit au seul élément tab[1]) et le
sous-tableau des autres valeurs non encore triées (de l’indice 2 à 20).
On recommence ce processus dans ce second sous-tableau : le minimum est
à présent l'élément d’indice 17 et de valeur 19.
Celui-ci viendra donc en 2\textsuperscript{ème} position, échangeant sa
place avec la valeur 52 qui s’y trouvait :}

\begin{center}
\tablehead{}
\begin{supertabular}{m{0.49700004cm}m{0.49700004cm}m{0.49700004cm}m{0.49700004cm}m{0.49700004cm}m{0.49700004cm}m{0.49700004cm}m{0.49700004cm}m{0.49700004cm}m{0.49700004cm}m{0.49700004cm}m{0.49700004cm}m{0.49700004cm}m{0.49700004cm}m{0.49700004cm}m{0.49700004cm}m{0.49700004cm}m{0.49700004cm}m{0.49700004cm}m{0.532cm}}
\itshape 1 &
\itshape 2 &
\itshape 3 &
\itshape 4 &
\itshape 5 &
\itshape 6 &
\itshape 7 &
\itshape 8 &
\itshape 9 &
\itshape 10 &
\itshape 11 &
\itshape 12 &
\itshape 13 &
\itshape 14 &
\itshape 15 &
\itshape 16 &
\itshape 17 &
\itshape 18 &
\itshape 19 &
\itshape 20\\\hline
\multicolumn{1}{|m{0.49700004cm}|}{ 17} &
\multicolumn{1}{m{0.49700004cm}|}{ 19} &
\multicolumn{1}{m{0.49700004cm}|}{ 61} &
\multicolumn{1}{m{0.49700004cm}|}{ 47} &
\multicolumn{1}{m{0.49700004cm}|}{ 82} &
\multicolumn{1}{m{0.49700004cm}|}{ 64} &
\multicolumn{1}{m{0.49700004cm}|}{ 95} &
\multicolumn{1}{m{0.49700004cm}|}{ 66} &
\multicolumn{1}{m{0.49700004cm}|}{ 84} &
\multicolumn{1}{m{0.49700004cm}|}{ 20} &
\multicolumn{1}{m{0.49700004cm}|}{ 32} &
\multicolumn{1}{m{0.49700004cm}|}{ 24} &
\multicolumn{1}{m{0.49700004cm}|}{ 46} &
\multicolumn{1}{m{0.49700004cm}|}{ 48} &
\multicolumn{1}{m{0.49700004cm}|}{ 75} &
\multicolumn{1}{m{0.49700004cm}|}{ 55} &
\multicolumn{1}{m{0.49700004cm}|}{ 52} &
\multicolumn{1}{m{0.49700004cm}|}{ 61} &
\multicolumn{1}{m{0.49700004cm}|}{ 21} &
\multicolumn{1}{m{0.532cm}|}{ 30}\\\hline
\end{supertabular}
\end{center}
{
Le sous-tableau trié est maintenant formé des deux premiers éléments, et
le sous-tableau non trié par les 18 éléments suivants. À
l’étape~\textit{i} de cet algorithme, on sélectionne donc le minimum du
sous-tableau non trié (entre les indices \textit{i} et \textit{n}). Une
fois le minimum localisé, on l’échange avec l’élément d’indice
\textit{i}. À chaque étape, la taille du sous-tableau trié augmente de
1 et celle du sous-tableau non trié diminue de 1. L’algorithme s’arrête
à la \textit{n}\textsuperscript{ème} étape, lorsque le sous-tableau
trié correspond au tableau de départ. En pratique l’arrêt se fait après
l’étape \textit{n} – 1, car lorsque le sous-tableau non trié n’est plus
que de taille 1, il contient nécessairement le maximum de l’ensemble.}

{\sffamily\bfseries\upshape
Algorithme}

{\sffamily
\textbf{module} triSélectionMinimaSuccessifs( tab↓↑ : tableau[1 à n]
d’entiers )}

{\sffamily
\ \ i, indiceMin : entiers}

{\sffamily
\ \ \textbf{pour} i \textbf{de} 1 \textstyleMotCl{à} n – 1
\textbf{faire}\ \ // i correspond à l’étape de l’algorithme}

{\sffamily
\ \ \ \ indiceMin \textrm{[F0DF?]} positionMin( tab, i, n )\ \ }

{\sffamily
\ \ \ \ {swap( tab[ i ], tab[ indiceMin ] )}}

{\sffamily
{\ \ }{\textbf{fin}}{
}{\textbf{pour}}}

{\sffamily
{\textbf{fin}}{
}{\textbf{module}}}


\bigskip

{\sffamily
// recherche de l’indice du minimum entre les indices début et fin }

{\sffamily
\textbf{module} positionMin(tab↓:tableau[1 à n] d’entiers,début,fin :
entiers ) → entier}

{\sffamily
\ \ indiceMin, i : entiers}

{\sffamily
\ \ indiceMin {\textrm{[F0DF?]}} début}

{\sffamily
\ \ \textbf{pour} i \textbf{de} début+1 \textstyleMotCl{à} fin
\textbf{faire}\ \ \ \ }

{\sffamily
\ \ \ \ \textstyleMotCl{si} tab[i] {\textless} tab[indiceMin]
\textstyleMotCl{alors}}

{\sffamily
\ \ \ \ \ \ indiceMin {\textrm{[F0DF?]}} i}

{\sffamily
\ \ \ \ \textstyleMotCl{fin} \textstyleMotCl{si}}

{\sffamily
{\ \ }{\textbf{fin}}{
}{\textbf{pour}}}

{\sffamily
\ \ \textstyleMotCl{retourner} indiceMin}

{\sffamily
{\textbf{fin}}{
}{\textbf{module}}}


\bigskip

{\sffamily
// échange entre 2 variables}

{\sffamily
\textbf{module} swap( a ↓↑, b ↓↑: entiers )}

{\sffamily
\ \ aux : entier}

{\sffamily
\ \ aux {\textrm{[F0DF?]}} a}

{\sffamily
\ \ a {\textrm{[F0DF?]}} b}

{\sffamily
\ \ b {\textrm{[F0DF?]}} aux}

{\sffamily
{\textbf{fin}}{
}{\textbf{module}}}


\bigskip

{\sffamily\bfseries\upshape
Exercices}

\liststyleExercice
\begin{enumerate}
\item {
Développez un algorithme similaire consistant à trier un tableau par
ordre croissant par sélection des \textbf{maxima} successifs. Le
sous-tableau trié apparaitra donc à droite du tableau, et les maxima
sélectionnés seront à chaque étape positionnés à droite du sous-tableau
non trié. }
\item {
Développez un algorithme combinant les deux recherches. À chaque étape,
on sélectionne donc le minimum et le maximum du sous-tableau restant à
trier et on les positionnera à l’endroit \textit{ad hoc}. N.B. : on
constatera que cette méthode n’apporte pas d’amélioration, ni en temps,
ni en simplicité, aux deux algorithmes de base !}
\end{enumerate}
\section{Tri bulle}
{
Il s’agit d’un tri par \textbf{permutations} ayant pour but d’amener à
chaque étape à la «~surface~» du sous-tableau non trié (on entend par
là l’élément d’indice minimum) la valeur la plus petite, appelée la
\textbf{bulle}. La caractéristique de cette méthode est que les
comparaisons ne se font qu’entre éléments consécutifs du tableau.}

{\sffamily\bfseries\upshape
Description de l’algorithme}

{
Prenons pour exemple un tableau de taille 14. En partant de la fin du
tableau, on le parcourt vers la gauche en comparant chaque couple de
valeurs consécutives. Quand deux valeurs sont dans le désordre, on les
permute. Le premier parcours s’achève lorsqu’on arrive à l’élément
d’indice 1 qui contient alors la «~bulle~»,
c'est-à-dire la plus petite valeur du tableau, soit 1
:}


\bigskip

\begin{center}
\tablehead{}
\begin{supertabular}{m{0.81100005cm}m{0.81100005cm}m{0.81100005cm}m{0.81100005cm}m{0.81100005cm}m{0.81100005cm}m{0.81100005cm}m{0.81100005cm}m{0.81100005cm}m{0.81100005cm}m{0.81100005cm}m{0.81100005cm}m{0.81100005cm}m{0.807cm}}
\centering \sffamily\itshape 1 &
\centering \sffamily\itshape 2 &
\centering \sffamily\itshape 3 &
\centering \sffamily\itshape 4 &
\centering \sffamily\itshape 5 &
\centering \sffamily\itshape 6 &
\centering \sffamily\itshape 7 &
\centering \sffamily\itshape 8 &
\centering \sffamily\itshape 9 &
\centering \sffamily\itshape 10 &
\centering \sffamily\itshape 11 &
\centering \sffamily\itshape 12 &
\centering \sffamily\itshape 13 &
\centering\arraybslash \sffamily\itshape
14\\\hline
\multicolumn{1}{|m{0.81100005cm}|}{\centering
\sffamily 10} &
\multicolumn{1}{m{0.81100005cm}|}{\centering
\sffamily 5} &
\multicolumn{1}{m{0.81100005cm}|}{\centering
\sffamily 12} &
\multicolumn{1}{m{0.81100005cm}|}{\centering
\sffamily 15} &
\multicolumn{1}{m{0.81100005cm}|}{\centering
\sffamily 4} &
\multicolumn{1}{m{0.81100005cm}|}{\centering
\sffamily 8} &
\multicolumn{1}{m{0.81100005cm}|}{\centering
\sffamily 1} &
\multicolumn{1}{m{0.81100005cm}|}{\centering
\sffamily 7} &
\multicolumn{1}{m{0.81100005cm}|}{\centering
\sffamily 12} &
\multicolumn{1}{m{0.81100005cm}|}{\centering
\sffamily 11} &
\multicolumn{1}{m{0.81100005cm}|}{\centering
\sffamily 3} &
\multicolumn{1}{m{0.81100005cm}|}{\centering
\sffamily 6} &
\multicolumn{1}{m{0.81100005cm}|}{\centering
\sffamily 5} &
\multicolumn{1}{m{0.807cm}|}{\centering\arraybslash
\sffamily 4}\\\hline
\end{supertabular}
\end{center}

\bigskip

\begin{flushleft}
\tablehead{}
\begin{supertabular}{|m{0.812cm}|m{0.812cm}|m{0.812cm}|m{0.812cm}|m{0.812cm}|m{0.812cm}|m{0.812cm}|m{0.812cm}|m{0.812cm}|m{0.812cm}|m{0.812cm}|m{0.812cm}|m{0.812cm}|m{0.828cm}|}
\hline
\centering \sffamily 10 &
\centering \sffamily 5 &
\centering \sffamily 12 &
\centering \sffamily 15 &
\centering \sffamily 4 &
\centering \sffamily 8 &
\centering \sffamily 1 &
\centering \sffamily 7 &
\centering \sffamily 12 &
\centering \sffamily 11 &
\centering \sffamily 3 &
\centering \sffamily 6 &
\centering \sffamily 4 &
\centering\arraybslash \sffamily 5\\\hline
\end{supertabular}
\end{flushleft}

\bigskip

\begin{flushleft}
\tablehead{}
\begin{supertabular}{|m{0.812cm}|m{0.812cm}|m{0.812cm}|m{0.812cm}|m{0.812cm}|m{0.812cm}|m{0.812cm}|m{0.812cm}|m{0.812cm}|m{0.812cm}|m{0.812cm}|m{0.812cm}|m{0.812cm}|m{0.828cm}|}
\hline
\centering \sffamily 10 &
\centering \sffamily 5 &
\centering \sffamily 12 &
\centering \sffamily 15 &
\centering \sffamily 4 &
\centering \sffamily 8 &
\centering \sffamily 1 &
\centering \sffamily 7 &
\centering \sffamily 12 &
\centering \sffamily 11 &
\centering \sffamily 3 &
\centering \sffamily 4 &
\centering \sffamily 6 &
\centering\arraybslash \sffamily 5\\\hline
\end{supertabular}
\end{flushleft}

\bigskip

\begin{flushleft}
\tablehead{}
\begin{supertabular}{|m{0.812cm}|m{0.812cm}|m{0.812cm}|m{0.812cm}|m{0.812cm}|m{0.812cm}|m{0.812cm}|m{0.812cm}|m{0.812cm}|m{0.812cm}|m{0.812cm}|m{0.812cm}|m{0.812cm}|m{0.828cm}|}
\hline
\centering \sffamily 10 &
\centering \sffamily 5 &
\centering \sffamily 12 &
\centering \sffamily 15 &
\centering \sffamily 4 &
\centering \sffamily 8 &
\centering \sffamily 1 &
\centering \sffamily 7 &
\centering \sffamily 12 &
\centering \sffamily 11 &
\centering \sffamily 3 &
\centering \sffamily 4 &
\centering \sffamily 6 &
\centering\arraybslash \sffamily 5\\\hline
\end{supertabular}
\end{flushleft}

\bigskip

\begin{flushleft}
\tablehead{}
\begin{supertabular}{|m{0.812cm}|m{0.812cm}|m{0.812cm}|m{0.812cm}|m{0.812cm}|m{0.812cm}|m{0.812cm}|m{0.812cm}|m{0.812cm}|m{0.812cm}|m{0.812cm}|m{0.812cm}|m{0.812cm}|m{0.828cm}|}
\hline
\centering \sffamily 10 &
\centering \sffamily 5 &
\centering \sffamily 12 &
\centering \sffamily 15 &
\centering \sffamily 4 &
\centering \sffamily 8 &
\centering \sffamily 1 &
\centering \sffamily 7 &
\centering \sffamily 12 &
\centering \sffamily 3 &
\centering \sffamily 11 &
\centering \sffamily 4 &
\centering \sffamily 6 &
\centering\arraybslash \sffamily 5\\\hline
\end{supertabular}
\end{flushleft}

\bigskip

\begin{flushleft}
\tablehead{}
\begin{supertabular}{|m{0.812cm}|m{0.812cm}|m{0.812cm}|m{0.812cm}|m{0.812cm}|m{0.812cm}|m{0.812cm}|m{0.812cm}|m{0.812cm}|m{0.812cm}|m{0.812cm}|m{0.812cm}|m{0.812cm}|m{0.828cm}|}
\hline
\centering \sffamily 10 &
\centering \sffamily 5 &
\centering \sffamily 12 &
\centering \sffamily 15 &
\centering \sffamily 4 &
\centering \sffamily 8 &
\centering \sffamily 1 &
\centering \sffamily 7 &
\centering \sffamily 3 &
\centering \sffamily 12 &
\centering \sffamily 11 &
\centering \sffamily 4 &
\centering \sffamily 6 &
\centering\arraybslash \sffamily 5\\\hline
\end{supertabular}
\end{flushleft}

\bigskip

\begin{flushleft}
\tablehead{}
\begin{supertabular}{|m{0.812cm}|m{0.812cm}|m{0.812cm}|m{0.812cm}|m{0.812cm}|m{0.812cm}|m{0.812cm}|m{0.812cm}|m{0.812cm}|m{0.812cm}|m{0.812cm}|m{0.812cm}|m{0.812cm}|m{0.828cm}|}
\hline
\centering \sffamily 10 &
\centering \sffamily 5 &
\centering \sffamily 12 &
\centering \sffamily 15 &
\centering \sffamily 4 &
\centering \sffamily 8 &
\centering \sffamily 1 &
\centering \sffamily 3 &
\centering \sffamily 7 &
\centering \sffamily 12 &
\centering \sffamily 11 &
\centering \sffamily 4 &
\centering \sffamily 6 &
\centering\arraybslash \sffamily 5\\\hline
\end{supertabular}
\end{flushleft}

\bigskip

\begin{flushleft}
\tablehead{}
\begin{supertabular}{|m{0.812cm}|m{0.812cm}|m{0.812cm}|m{0.812cm}|m{0.812cm}|m{0.812cm}|m{0.812cm}|m{0.812cm}|m{0.812cm}|m{0.812cm}|m{0.812cm}|m{0.812cm}|m{0.812cm}|m{0.828cm}|}
\hline
\centering \sffamily 10 &
\centering \sffamily 5 &
\centering \sffamily 12 &
\centering \sffamily 15 &
\centering \sffamily 4 &
\centering \sffamily 8 &
\centering \sffamily 1 &
\centering \sffamily 3 &
\centering \sffamily 7 &
\centering \sffamily 12 &
\centering \sffamily 11 &
\centering \sffamily 4 &
\centering \sffamily 6 &
\centering\arraybslash \sffamily 5\\\hline
\end{supertabular}
\end{flushleft}

\bigskip

\begin{flushleft}
\tablehead{}
\begin{supertabular}{|m{0.812cm}|m{0.812cm}|m{0.812cm}|m{0.812cm}|m{0.812cm}|m{0.812cm}|m{0.812cm}|m{0.812cm}|m{0.812cm}|m{0.812cm}|m{0.812cm}|m{0.812cm}|m{0.812cm}|m{0.828cm}|}
\hline
\centering \sffamily 10 &
\centering \sffamily 5 &
\centering \sffamily 12 &
\centering \sffamily 15 &
\centering \sffamily 4 &
\centering \sffamily 1 &
\centering \sffamily 8 &
\centering \sffamily 3 &
\centering \sffamily 7 &
\centering \sffamily 12 &
\centering \sffamily 11 &
\centering \sffamily 4 &
\centering \sffamily 6 &
\centering\arraybslash \sffamily 5\\\hline
\end{supertabular}
\end{flushleft}

\bigskip

\begin{flushleft}
\tablehead{}
\begin{supertabular}{|m{0.812cm}|m{0.812cm}|m{0.812cm}|m{0.812cm}|m{0.812cm}|m{0.812cm}|m{0.812cm}|m{0.812cm}|m{0.812cm}|m{0.812cm}|m{0.812cm}|m{0.812cm}|m{0.812cm}|m{0.828cm}|}
\hline
\centering \sffamily 10 &
\centering \sffamily 5 &
\centering \sffamily 12 &
\centering \sffamily 15 &
\centering \sffamily 1 &
\centering \sffamily 4 &
\centering \sffamily 8 &
\centering \sffamily 3 &
\centering \sffamily 7 &
\centering \sffamily 12 &
\centering \sffamily 11 &
\centering \sffamily 4 &
\centering \sffamily 6 &
\centering\arraybslash \sffamily 5\\\hline
\end{supertabular}
\end{flushleft}

\bigskip

\begin{flushleft}
\tablehead{}
\begin{supertabular}{|m{0.812cm}|m{0.812cm}|m{0.812cm}|m{0.812cm}|m{0.812cm}|m{0.812cm}|m{0.812cm}|m{0.812cm}|m{0.812cm}|m{0.812cm}|m{0.812cm}|m{0.812cm}|m{0.812cm}|m{0.828cm}|}
\hline
\centering \sffamily 10 &
\centering \sffamily 5 &
\centering \sffamily 12 &
\centering \sffamily 1 &
\centering \sffamily 15 &
\centering \sffamily 4 &
\centering \sffamily 8 &
\centering \sffamily 3 &
\centering \sffamily 7 &
\centering \sffamily 12 &
\centering \sffamily 11 &
\centering \sffamily 4 &
\centering \sffamily 6 &
\centering\arraybslash \sffamily 5\\\hline
\end{supertabular}
\end{flushleft}

\bigskip

\begin{flushleft}
\tablehead{}
\begin{supertabular}{|m{0.812cm}|m{0.812cm}|m{0.812cm}|m{0.812cm}|m{0.812cm}|m{0.812cm}|m{0.812cm}|m{0.812cm}|m{0.812cm}|m{0.812cm}|m{0.812cm}|m{0.812cm}|m{0.812cm}|m{0.828cm}|}
\hline
\centering \sffamily 10 &
\centering \sffamily 5 &
\centering \sffamily 1 &
\centering \sffamily 12 &
\centering \sffamily 15 &
\centering \sffamily 4 &
\centering \sffamily 8 &
\centering \sffamily 3 &
\centering \sffamily 7 &
\centering \sffamily 12 &
\centering \sffamily 11 &
\centering \sffamily 4 &
\centering \sffamily 6 &
\centering\arraybslash \sffamily 5\\\hline
\end{supertabular}
\end{flushleft}

\bigskip

\begin{flushleft}
\tablehead{}
\begin{supertabular}{|m{0.812cm}|m{0.812cm}|m{0.812cm}|m{0.812cm}|m{0.812cm}|m{0.812cm}|m{0.812cm}|m{0.812cm}|m{0.812cm}|m{0.812cm}|m{0.812cm}|m{0.812cm}|m{0.812cm}|m{0.828cm}|}
\hline
\centering \sffamily 10 &
\centering \sffamily 1 &
\centering \sffamily 5 &
\centering \sffamily 12 &
\centering \sffamily 15 &
\centering \sffamily 4 &
\centering \sffamily 8 &
\centering \sffamily 3 &
\centering \sffamily 7 &
\centering \sffamily 12 &
\centering \sffamily 11 &
\centering \sffamily 4 &
\centering \sffamily 6 &
\centering\arraybslash \sffamily 5\\\hline
\end{supertabular}
\end{flushleft}

\bigskip

\begin{flushleft}
\tablehead{}
\begin{supertabular}{|m{0.812cm}|m{0.812cm}|m{0.812cm}|m{0.812cm}|m{0.812cm}|m{0.812cm}|m{0.812cm}|m{0.812cm}|m{0.812cm}|m{0.812cm}|m{0.812cm}|m{0.812cm}|m{0.812cm}|m{0.828cm}|}
\hline
\centering \sffamily 1 &
\centering \sffamily 10 &
\centering \sffamily 5 &
\centering \sffamily 12 &
\centering \sffamily 15 &
\centering \sffamily 4 &
\centering \sffamily 8 &
\centering \sffamily 3 &
\centering \sffamily 7 &
\centering \sffamily 12 &
\centering \sffamily 11 &
\centering \sffamily 4 &
\centering \sffamily 6 &
\centering\arraybslash \sffamily 5\\\hline
\end{supertabular}
\end{flushleft}

\bigskip

{
Ceci achève le premier parcours. On recommence à présent le même
processus dans le sous-tableau débutant au deuxième élément, ce qui
donnera le contenu suivant :}


\bigskip

\begin{flushleft}
\tablehead{}
\begin{supertabular}{|m{0.812cm}|m{0.812cm}|m{0.812cm}|m{0.812cm}|m{0.812cm}|m{0.812cm}|m{0.812cm}|m{0.812cm}|m{0.812cm}|m{0.812cm}|m{0.812cm}|m{0.812cm}|m{0.812cm}|m{0.828cm}|}
\hline
\centering \sffamily 1 &
\centering \sffamily 3 &
\centering \sffamily 10 &
\centering \sffamily 5 &
\centering \sffamily 12 &
\centering \sffamily 15 &
\centering \sffamily 4 &
\centering \sffamily 8 &
\centering \sffamily 4 &
\centering \sffamily 7 &
\centering \sffamily 12 &
\centering \sffamily 11 &
\centering \sffamily 5 &
\centering\arraybslash \sffamily 6\\\hline
\end{supertabular}
\end{flushleft}
{
Comme pour le tri par recherche des minima successifs, à l’issue de
l’étape \textit{i} de l’algorithme, on obtient un sous-tableau trié
(entre les indices 1 et \textit{i }) et un tableau non encore trié
(entre les indices \textit{i} + 1 et \textit{n}). Le tri est donc
achevé à l’issue de l’étape \textit{n} – 1. Cependant, on notera que
petit à petit, outre le déplacement de la bulle, les éléments plus
petits évoluent progressivement vers la gauche, et les plus gros vers
la droite, de sorte qu’il est fréquent que le nombre d’étapes effectif
est inférieur au nombre d’étapes théorique.}

{\sffamily\bfseries\upshape
Algorithme}

{
Dans cet algorithme, la variable \textstyleCodeInsr{indiceBulle} est à
la fois le compteur d’étapes et aussi l’indice de l’élément récepteur
de la bulle à la fin d’une étape donnée.}


\bigskip

{\sffamily
\textbf{module} triBulle( tab↓↑: tableau[1 à n] d’entiers )}

{\sffamily
\ \ indiceBulle, j : entiers}

{\sffamily
\ \ \textbf{pour} indiceBulle \textstyleMotCl{de} 1 \textstyleMotCl{à} n
– 1 \textbf{faire}}

{\sffamily
\ \ \ \ \textbf{pour} j \textstyleMotCl{de} n – 1 \textstyleMotCl{à}
indiceBulle \textstyleMotCl{par} – 1 \textbf{faire}}

{\sffamily
\ \ \ \ \ \ \textstyleMotCl{si} tab[j] {\textgreater} tab[j + 1]
\textstyleMotCl{alors}\ \ \ \ \ \ \ \ \ \ }

{\sffamily
\ \ \ \ \ \ \ \ swap( tab[ j ], tab[ j+1 ] ) \ // voir algorithme
précédent}

{\sffamily
\ \ \ \ \ \ \textstyleMotCl{fin} \textstyleMotCl{si}}

{\sffamily
\ \ \ \ {\textbf{fin}}{
}{\textbf{pour}}}

{\sffamily
\ \ {\textbf{fin}}{
}{\textbf{pour}}}

{\sffamily
{\textbf{fin}}{
}{\textbf{module}}}


\bigskip

{\sffamily\bfseries\upshape
Exercices}

\liststyleExercice
\setcounter{saveenum}{\value{enumi}}
\begin{enumerate}
\setcounter{enumi}{\value{saveenum}}
\item {
Écrire une amélioration du tri bulle consistant à mémoriser à chaque
étape l’indice de la dernière permutation, celui-ci délimitant en fait
la véritable taille du sous-tableau trié à l’issue d’une étape. En lieu
et place de la boucle \textstyleMotCl{pour} n’incrémentant l’indice
bulle que de 1 à la fois, on écrira une boucle \textstyleMotCl{tant
que} à l’issue de laquelle \textstyleCodeInsr{indiceBulle} prendra la
valeur de l’indice de dernière permutation + 1.}
\item {
L’amélioration précédente est issue de l’observation du sous-tableau
déjà trié en début du tableau initial. On peut de même étudier la
possibilité d’avoir un sous-tableau, trié également, mais cadré à
droite dans le tableau à trier. Cette symétrie suggère une amélioration
supplémentaire qui consiste à changer de sens à la fin d’un parcours
pour entamer le parcours suivant. Lorsqu’on change de sens, on amènera
l’élément le plus grand (qu’on peut nommer le «~plomb~») au fond du
tableau non trié. L’association des deux méthodes donne ce qu’on
appelle le \textbf{tri shaker}, dont le but est de restreindre le
sous-tableau non trié en augmentant sa borne inférieure et en diminuant
sa borne supérieure. Écrire l’algorithme qui réalise cette méthode de
tri.}
\item {
Écrire un algorithme qui trie un tableau de (variables structurées)
dates par ordre chronologique croissant.}
\item {
Écrire un algorithme qui trie un tableau de personnes sur leur date
d'anniversaire (sans tenir compte de
l'année). Deux personnes nées le même jour seront
départagées par leur nom.}

{
Une personne est représentée par la structure }
\end{enumerate}
{\sffamily
\textstyleMotCl{structure} Personne \textstyleMotCl{composée}
\textstyleMotCl{de}}

{\sffamily
\ \ nom : chaine}

{\sffamily
\ \ dateNaissance : Date}

{\sffamily
\textstyleMotCl{fin} \textstyleMotCl{structure}}


\bigskip

\section{Cas particuliers}
{\sffamily\bfseries\upshape
Tri partiel}

{
Parfois, on n’est pas intéressé par un tri complet de l’ensemble mais on
veut uniquement les «~k~» plus petites (ou plus grandes) valeurs. Le
tri par recherche des minima et le tri bulle sont particulièrement
adaptés à cette situation ; il suffit de les arrêter plus tôt. Par
contre, le tri par insertion est inefficace car il ne permet pas un
arrêt anticipé.}

{\sffamily\bfseries\upshape
Ensemble (presque) trié}

{
Les trois tris que nous avons vus ont la même complexité quadratique
(voir plus loin) dans le cas général mais comment se comportent-ils
dans le cas d’un tableau (presque) trié ? Dans ce cas, on peut vérifier
que le tri par recherche des minima successifs reste quadratique alors
que les deux autres deviennent linéaires (en n).}

{\sffamily\bfseries\upshape
Sélection des meilleurs}

{
Une autre situation courante est la suivante : \ un ensemble de valeurs
sont traitées (lues, calculées, ...) et il ne faut \ garder que les
\textstyleCodeInsr{k} plus petites (ou plus grandes).
L'algorithme est plus simple à écrire si on utilise
une position supplémentaire en fin de tableau. Notons aussi qu’il faut
tenir compte du cas où il y a moins de valeurs que le nombre de valeurs
voulues ; c’est pourquoi on ajoute une variable indiquant le nombre
exact de valeurs dans le tableau.}

{\sffamily\bfseries\scshape
Exemple : les k plus petites valeurs}

{
On lit un ensemble de valeurs strictement positives (un 0 indique la fin
de la lecture) et on ne garde que les «~k~» plus petites valeurs.}


\bigskip


\bigskip

{\sffamily
\textstyleMotCl{module} meilleurs ( plusPetits~↓↑ :
\textstyleMotCl{tableau} [ 1 à k+1 ] d’entiers, }

{\sffamily
\ \ \ \ \ \ \ \ \ \ \ \ \ \ \ \ \ \ nbValeurs~↑ : entier )}

{\sffamily
\ \ val : entier}

{\sffamily
\ \ nbValeurs \textrm{[F0DF?]} 0}

{\sffamily
\ \ \textstyleMotCl{lire} val}

{\sffamily
\ \ \textstyleMotCl{tant} \textstyleMotCl{que} val ${\neq}$ 0
\textstyleMotCl{faire}}

{\sffamily
\ \ \ \ insérer( val, plusPetits, nbValeurs )}

{\sffamily
\ \ \ \ \textstyleMotCl{lire} val}

{\sffamily
\ \ \textstyleMotCl{fin} \textstyleMotCl{tant} \textstyleMotCl{que}}

{\sffamily
\textstyleMotCl{fin} \textstyleMotCl{module}}


\bigskip

{\sffamily
\textstyleMotCl{module} insérer ( val: entier, plusPetits~↓↑:
\textstyleMotCl{tableau} [ 1 à k+1 ] d’entiers, }

{\sffamily
\ \ \ \ \ \ \ \ \ \ \ \ \ \ \ \ \ \ nbValeurs~↓↑: entier )}

{\sffamily
\ \ i : entier}

{\sffamily
\ \ i \textrm{[F0DF?]} nbValeurs}

{\sffamily
\ \ \textstyleMotCl{tant} \textstyleMotCl{que} i {\textgreater} 0 ET val
{\textless} plusPetits~[ i ] \textstyleMotCl{faire}}

{\sffamily
\ \ \ \ plusPetits~{[ i + 1 ]
}\textrm{[F0DF?]}{ plusPetits~[ i ]}}

{\sffamily
{\ \ \ \ }i \textrm{[F0DF?]} i - 1}

{\sffamily
\ \ \textstyleMotCl{fin} \textstyleMotCl{tant} \textstyleMotCl{que}}

{\sffamily
\ \ plusPetits~{[ i + 1 ]
}\textrm{[F0DF?]}{ val}}

{\sffamily
\ \ \textstyleMotCl{si} nbValeurs {\textless} k
\textstyleMotCl{alors}\ \ \ \ \ \ \ \ \ \ }

{\sffamily
\ \ \ \ nbValeurs{
}\textrm{[F0DF?]}{ }nbValeurs + 1}

{\sffamily
\ \ \textstyleMotCl{fin} \textstyleMotCl{si}}

{\sffamily
\textstyleMotCl{fin} \textstyleMotCl{module}}

\section{Recherche dichotomique}
{
La recherche dichotomique (ou binaire) a pour essence de réduire à
chaque étape la taille de l’ensemble de recherche de moitié, jusqu’à ce
qu’il ne reste qu’un seul élément dont la valeur devrait être celle
recherchée, sauf bien entendu en cas d’inexistence de cette valeur dans
l’ensemble de départ. Nous l’expliquons pour un tableau ordonné mais
elle peut s’appliquer aussi avec des changements mineurs pour
d'autres structures (par exemple la liste ordonnée que
nous étudierons plus tard).}

{\sffamily\bfseries\upshape
Description de l’algorithme}

{
Soit \textstyleCodeInsr{val} la valeur recherchée dans une zone
délimitée par les indices i\textstyleCodeInsr{ndiceGauche} et
i\textstyleCodeInsr{ndiceDroit}. On commence par déterminer l’élément
médian, c’est-à-dire celui qui se trouve « au milieu » de la zone de
recherche ; son indice sera déterminé par la formule }

{\centering
\textstyleCodeInsr{{indiceMédian}}{
}\textsf{←}{
}\textstyleCodeInsr{{(indiceGauche +
indiceDroit)
}}\textstyleCodeInsr{{DIV}}\textstyleCodeInsr{{
2}}
\par}

{
{On compare alors
}\textstyleCodeInsr{{val}}{
avec la valeur de cet élément médian ; il est possible qu’on ait trouvé
la valeur cherchée; sinon, on partage la zone de recherche en deux
parties :}\textstyleWWPolicepardfaut{ une qui
}\textstyleWWPolicepardfaut{\textbf{ne contient certainement
pas}}\textstyleWWPolicepardfaut{ la valeur cherchée et une qui
}\textstyleWWPolicepardfaut{\textbf{pourrait la
contenir}}\textstyleWWPolicepardfaut{. C’est cette deuxième partie qui
devient}{ la nouvelle zone de recherche. On
réitère le processus jusqu’à ce que la valeur cherchée soit trouvée ou
que la zone de recherche soit réduite à sa plus simple expression,
c’est-à-dire un seul élément.}}

{\sffamily\bfseries
Algorithme}

{\sffamily
\textstyleMotCl{module} rechercheDichotomique( tab~↓↑ : tableau [1 à n]
de T,
\ \ \ \ \ \ \ \ \ \ \ \ \ \ valeur~↓ : T, pos~↑ : entier ) → booléen}

{\sffamily

\ \ \textstyleWWPolicepardfaut{indiceDroit, indiceGauche, indiceMédian :
entiers}}

{\sffamily
\textstyleWWPolicepardfaut{\ \ candidat : T}}

{\sffamily
\textstyleWWPolicepardfaut{\ \ trouvé : booléen}}


\bigskip

{\sffamily
\ \ \textstyleWWPolicepardfaut{indiceGauche}\textrm{ }←
\textstyleWWPolicepardfaut{1}}

{\sffamily
\ \ \textstyleWWPolicepardfaut{indiceDroit}\textrm{ }←
\textstyleWWPolicepardfaut{n}}

{\sffamily
\ \ \textstyleWWPolicepardfaut{trouvé}\textrm{ }←
\textstyleWWPolicepardfaut{faux}}


\bigskip

{\sffamily
\textstyleWWPolicepardfaut{\ \ }\textstyleMotCl{tant
que}\textstyleWWPolicepardfaut{
}\textstyleWWPolicepardfaut{NON}\textstyleWWPolicepardfaut{ trouvé
}\textstyleWWPolicepardfaut{ET}\textstyleWWPolicepardfaut{ indiceGauche
}\textstyleWWPolicepardfaut{${\leq}$}\textstyleWWPolicepardfaut{
indiceDroit \ }\textstyleMotCl{faire}}

{\sffamily
\ \ \ \ \textstyleWWPolicepardfaut{indiceMédian }←
\textstyleWWPolicepardfaut{( indiceGauche + indiceDroit )
}\textstyleWWPolicepardfaut{DIV}\textstyleWWPolicepardfaut{ 2}}

{\sffamily
\ \ \ \ \textstyleWWPolicepardfaut{candidat}\textrm{ }←
tab[\textstyleWWPolicepardfaut{ indiceMédian ]}}

{\sffamily
\textstyleWWPolicepardfaut{\ \ \ \ }\textstyleMotCl{selon que}}

{\sffamily
\ \ \ \  \ \ \ \textstyleWWPolicepardfaut{candidat~ = valeur :}}

{\sffamily
\ \ \ \ \ \ \textstyleWWPolicepardfaut{trouvé}\textrm{ }←
\textstyleWWPolicepardfaut{vrai}}

{\sffamily
\textstyleWWPolicepardfaut{\ \ \ \  \ \ \ candidat {\textless} valeur
:}}

{\sffamily
\ \ \ \ \ \ \textstyleWWPolicepardfaut{indiceGauche}\textrm{ }←
\textstyleWWPolicepardfaut{indiceMédian + 1\ \ // on garde la partie
droite}}

{\sffamily
\textstyleWWPolicepardfaut{\ \ \ \  \ \ \ candidat {\textgreater} valeur
:}}

{\sffamily
\ \ \ \ \ \ \textstyleWWPolicepardfaut{indiceDroit}\textrm{ }←
\textstyleWWPolicepardfaut{indiceMédian~– 1\ \ // on garde la partie
gauche}}

{\sffamily
\textstyleWWPolicepardfaut{\ \ \ \ }\textstyleMotCl{fin}\textstyleWWPolicepardfaut{
}\textstyleMotCl{selon que}\textstyleWWPolicepardfaut{\ \ }}

{\sffamily
\textstyleWWPolicepardfaut{\ \ }\textstyleMotCl{fin}\textstyleWWPolicepardfaut{
}\textstyleMotCl{tant que}}

{\sffamily
\textstyleWWPolicepardfaut{\ \ }\textstyleMotCl{si}\textstyleWWPolicepardfaut{
trouvé }\textstyleMotCl{alors}}

{\sffamily
\ \ \ \ \textstyleWWPolicepardfaut{pos}\textrm{ }←
\textstyleWWPolicepardfaut{indiceMédian}}

{\sffamily
\textstyleWWPolicepardfaut{\ \ }\textstyleMotCl{sinon}\textstyleWWPolicepardfaut{\ \ \ \ \ \ }}

{\sffamily
\ \ \ \ \textstyleWWPolicepardfaut{pos}\textrm{ }←
\textstyleWWPolicepardfaut{indiceGauche\ \ // dans le cas où la valeur
n’est pas trouvée, }}

{\sffamily
\textstyleWWPolicepardfaut{\ \ \ \ \ \ \ \ \ \ \ \ // on vérifiera que
indiceGauche donne la }}

{\sffamily
\textstyleWWPolicepardfaut{\ \ \ \ \ \ \ \ \ \ \ \ // valeur où elle
pourrait être insérée.}}

{\sffamily
\textstyleWWPolicepardfaut{\ \ }\textstyleMotCl{fin}\textstyleWWPolicepardfaut{
}\textstyleMotCl{si}\textstyleWWPolicepardfaut{\ \ }}

{\sffamily
\textstyleWWPolicepardfaut{\ \ }\textstyleMotCl{retourner}\textstyleWWPolicepardfaut{
trouvé\ \ \ \ }}

{\sffamily
\textstyleMotCl{fin} \textstyleMotCl{module}}

{\sffamily\bfseries\upshape
Exemple de recherche fructueuse}

{
{Supposons que l’on recherche la valeur
}{\textbf{23}}{ dans un
tableau de 20 entiers. La zone ombrée représente à chaque fois la
partie de recherche, qui est au départ le tableau dans son entièreté.
Au départ,
}\textstyleCodeInsr{{indiceGauche}}{
= 1 et
}\textstyleCodeInsr{{indiceDroit}}{
= 20.}}

\begin{flushleft}
\tablehead{}
\begin{supertabular}{m{0.507cm}m{0.509cm}m{0.509cm}m{0.509cm}m{0.507cm}m{0.509cm}m{0.509cm}m{0.509cm}m{0.509cm}m{0.507cm}m{0.509cm}m{0.509cm}m{0.509cm}m{0.509cm}m{0.51100004cm}m{0.51100004cm}m{0.509cm}m{0.51100004cm}m{0.51100004cm}m{0.527cm}}
\centering \itshape 1 &
\centering \itshape 2 &
\centering \itshape 3 &
\centering \itshape 4 &
\centering \itshape 5 &
\centering \itshape 6 &
\centering \itshape 7 &
\centering \itshape 8 &
\centering \itshape 9 &
\centering \itshape 10 &
\centering \itshape 11 &
\centering \itshape 12 &
\centering \itshape 13 &
\centering \itshape 14 &
\centering \itshape 15 &
\centering \itshape 16 &
\centering \itshape 17 &
\centering \itshape 18 &
\centering \itshape 19 &
\centering\arraybslash \itshape 20\\\hline
\multicolumn{1}{|m{0.507cm}|}{ 1} &
\multicolumn{1}{m{0.509cm}|}{ 3} &
\multicolumn{1}{m{0.509cm}|}{ 5} &
\multicolumn{1}{m{0.509cm}|}{ 7} &
\multicolumn{1}{m{0.507cm}|}{ 7} &
\multicolumn{1}{m{0.509cm}|}{ 9} &
\multicolumn{1}{m{0.509cm}|}{ 9} &
\multicolumn{1}{m{0.509cm}|}{ 10} &
\multicolumn{1}{m{0.509cm}|}{ 10} &
\multicolumn{1}{m{0.507cm}|}{ 15} &
\multicolumn{1}{m{0.509cm}|}{ 16} &
\multicolumn{1}{m{0.509cm}|}{ 20} &
\multicolumn{1}{m{0.509cm}|}{ 23} &
\multicolumn{1}{m{0.509cm}|}{ 23} &
\multicolumn{1}{m{0.51100004cm}|}{ 25} &
\multicolumn{1}{m{0.51100004cm}|}{ 28} &
\multicolumn{1}{m{0.509cm}|}{ 28} &
\multicolumn{1}{m{0.51100004cm}|}{ 28} &
\multicolumn{1}{m{0.51100004cm}|}{ 29} &
\multicolumn{1}{m{0.527cm}|}{ 29}\\\hline
\end{supertabular}
\end{flushleft}

\bigskip

{
\textstyleWWPolicepardfaut{\textit{Première
étape}}\textstyleWWPolicepardfaut{ :
}\textstyleCodeInsr{indiceMédian}\textstyleWWPolicepardfaut{
}\textsf{←}\textstyleWWPolicepardfaut{ }\textstyleCodeInsr{(1 + 20)
}\textstyleCodeInsr{DIV}\textstyleCodeInsr{
2}\textstyleWWPolicepardfaut{, c’est-à-dire 10. Comme la valeur en
position 10 est 15, la valeur cherchée doit se trouver à sa droite, et
on prend comme nouvelle zone de recherche celle délimitée par
}\textstyleCodeInsr{indiceGauche}\textstyleWWPolicepardfaut{ = 11 et
}\textstyleCodeInsr{indiceDroit}\textstyleWWPolicepardfaut{ = 20.}}

\begin{flushleft}
\tablehead{}
\begin{supertabular}{m{0.507cm}m{0.509cm}m{0.509cm}m{0.509cm}m{0.507cm}m{0.509cm}m{0.509cm}m{0.509cm}m{0.509cm}m{0.507cm}m{0.509cm}m{0.509cm}m{0.509cm}m{0.509cm}m{0.51100004cm}m{0.51100004cm}m{0.509cm}m{0.51100004cm}m{0.51100004cm}m{0.527cm}}
\centering \itshape 1 &
\centering \itshape 2 &
\centering \itshape 3 &
\centering \itshape 4 &
\centering \itshape 5 &
\centering \itshape 6 &
\centering \itshape 7 &
\centering \itshape 8 &
\centering \itshape 9 &
\centering \itshape 10 &
\centering \itshape 11 &
\centering \itshape 12 &
\centering \itshape 13 &
\centering \itshape 14 &
\centering \itshape 15 &
\centering \itshape 16 &
\centering \itshape 17 &
\centering \itshape 18 &
\centering \itshape 19 &
\centering\arraybslash \itshape 20\\\hline
\multicolumn{1}{|m{0.507cm}|}{ 1} &
\multicolumn{1}{m{0.509cm}|}{ 3} &
\multicolumn{1}{m{0.509cm}|}{ 5} &
\multicolumn{1}{m{0.509cm}|}{ 7} &
\multicolumn{1}{m{0.507cm}|}{ 7} &
\multicolumn{1}{m{0.509cm}|}{ 9} &
\multicolumn{1}{m{0.509cm}|}{ 9} &
\multicolumn{1}{m{0.509cm}|}{ 10} &
\multicolumn{1}{m{0.509cm}|}{ 10} &
\multicolumn{1}{m{0.507cm}|}{ 15} &
\multicolumn{1}{m{0.509cm}|}{ 16} &
\multicolumn{1}{m{0.509cm}|}{ 20} &
\multicolumn{1}{m{0.509cm}|}{ 23} &
\multicolumn{1}{m{0.509cm}|}{ 23} &
\multicolumn{1}{m{0.51100004cm}|}{ 25} &
\multicolumn{1}{m{0.51100004cm}|}{ 28} &
\multicolumn{1}{m{0.509cm}|}{ 28} &
\multicolumn{1}{m{0.51100004cm}|}{ 28} &
\multicolumn{1}{m{0.51100004cm}|}{ 29} &
\multicolumn{1}{m{0.527cm}|}{ 29}\\\hline
\end{supertabular}
\end{flushleft}

\bigskip

{
\textstyleWWPolicepardfaut{\textit{Deuxième
étape}}\textstyleWWPolicepardfaut{ :
}\textstyleCodeInsr{indiceMédian}\textstyleWWPolicepardfaut{
}\textsf{←}\textstyleWWPolicepardfaut{ }\textstyleCodeInsr{(11 + 20)
}\textstyleCodeInsr{DIV}\textstyleCodeInsr{
2}\textstyleWWPolicepardfaut{, c’est-à-dire 15. Comme on y trouve la
valeur 25, on garde les éléments situées à la gauche de celui-ci ; la
nouvelle zone de recherche est [11, 14].}}

\begin{flushleft}
\tablehead{}
\begin{supertabular}{m{0.507cm}m{0.509cm}m{0.509cm}m{0.509cm}m{0.507cm}m{0.509cm}m{0.509cm}m{0.509cm}m{0.509cm}m{0.507cm}m{0.509cm}m{0.509cm}m{0.509cm}m{0.509cm}m{0.51100004cm}m{0.51100004cm}m{0.509cm}m{0.51100004cm}m{0.51100004cm}m{0.527cm}}
\centering \itshape 1 &
\centering \itshape 2 &
\centering \itshape 3 &
\centering \itshape 4 &
\centering \itshape 5 &
\centering \itshape 6 &
\centering \itshape 7 &
\centering \itshape 8 &
\centering \itshape 9 &
\centering \itshape 10 &
\centering \itshape 11 &
\centering \itshape 12 &
\centering \itshape 13 &
\centering \itshape 14 &
\centering \itshape 15 &
\centering \itshape 16 &
\centering \itshape 17 &
\centering \itshape 18 &
\centering \itshape 19 &
\centering\arraybslash \itshape 20\\\hline
\multicolumn{1}{|m{0.507cm}|}{ 1} &
\multicolumn{1}{m{0.509cm}|}{ 3} &
\multicolumn{1}{m{0.509cm}|}{ 5} &
\multicolumn{1}{m{0.509cm}|}{ 7} &
\multicolumn{1}{m{0.507cm}|}{ 7} &
\multicolumn{1}{m{0.509cm}|}{ 9} &
\multicolumn{1}{m{0.509cm}|}{ 9} &
\multicolumn{1}{m{0.509cm}|}{ 10} &
\multicolumn{1}{m{0.509cm}|}{ 10} &
\multicolumn{1}{m{0.507cm}|}{ 15} &
\multicolumn{1}{m{0.509cm}|}{ 16} &
\multicolumn{1}{m{0.509cm}|}{ 20} &
\multicolumn{1}{m{0.509cm}|}{ 23} &
\multicolumn{1}{m{0.509cm}|}{ 23} &
\multicolumn{1}{m{0.51100004cm}|}{ 25} &
\multicolumn{1}{m{0.51100004cm}|}{ 28} &
\multicolumn{1}{m{0.509cm}|}{ 28} &
\multicolumn{1}{m{0.51100004cm}|}{ 28} &
\multicolumn{1}{m{0.51100004cm}|}{ 29} &
\multicolumn{1}{m{0.527cm}|}{ 29}\\\hline
\end{supertabular}
\end{flushleft}

\bigskip

{
\textstyleWWPolicepardfaut{\textit{Troisième
étape}}\textstyleWWPolicepardfaut{ :
}\textstyleCodeInsr{indiceMédian}\textstyleWWPolicepardfaut{
}\textsf{←}\textstyleWWPolicepardfaut{ }\textstyleCodeInsr{(11 + 14)
}\textstyleCodeInsr{DIV}\textstyleCodeInsr{
2}\textstyleWWPolicepardfaut{, c’est-à-dire 12 où se trouve l’élément
20. La zone de recherche devient
}\textstyleCodeInsr{indiceGauche}\textstyleWWPolicepardfaut{ = 13 et
}\textstyleCodeInsr{indiceDroit}\textstyleWWPolicepardfaut{ = 14.}}

\begin{flushleft}
\tablehead{}
\begin{supertabular}{m{0.507cm}m{0.509cm}m{0.509cm}m{0.509cm}m{0.507cm}m{0.509cm}m{0.509cm}m{0.509cm}m{0.509cm}m{0.507cm}m{0.509cm}m{0.509cm}m{0.509cm}m{0.509cm}m{0.51100004cm}m{0.51100004cm}m{0.509cm}m{0.51100004cm}m{0.51100004cm}m{0.527cm}}
\centering \itshape 1 &
\centering \itshape 2 &
\centering \itshape 3 &
\centering \itshape 4 &
\centering \itshape 5 &
\centering \itshape 6 &
\centering \itshape 7 &
\centering \itshape 8 &
\centering \itshape 9 &
\centering \itshape 10 &
\centering \itshape 11 &
\centering \itshape 12 &
\centering \itshape 13 &
\centering \itshape 14 &
\centering \itshape 15 &
\centering \itshape 16 &
\centering \itshape 17 &
\centering \itshape 18 &
\centering \itshape 19 &
\centering\arraybslash \itshape 20\\\hline
\multicolumn{1}{|m{0.507cm}|}{ 1} &
\multicolumn{1}{m{0.509cm}|}{ 3} &
\multicolumn{1}{m{0.509cm}|}{ 5} &
\multicolumn{1}{m{0.509cm}|}{ 7} &
\multicolumn{1}{m{0.507cm}|}{ 7} &
\multicolumn{1}{m{0.509cm}|}{ 9} &
\multicolumn{1}{m{0.509cm}|}{ 9} &
\multicolumn{1}{m{0.509cm}|}{ 10} &
\multicolumn{1}{m{0.509cm}|}{ 10} &
\multicolumn{1}{m{0.507cm}|}{ 15} &
\multicolumn{1}{m{0.509cm}|}{ 16} &
\multicolumn{1}{m{0.509cm}|}{ 20} &
\multicolumn{1}{m{0.509cm}|}{ 23} &
\multicolumn{1}{m{0.509cm}|}{ 23} &
\multicolumn{1}{m{0.51100004cm}|}{ 25} &
\multicolumn{1}{m{0.51100004cm}|}{ 28} &
\multicolumn{1}{m{0.509cm}|}{ 28} &
\multicolumn{1}{m{0.51100004cm}|}{ 28} &
\multicolumn{1}{m{0.51100004cm}|}{ 29} &
\multicolumn{1}{m{0.527cm}|}{ 29}\\\hline
\end{supertabular}
\end{flushleft}

\bigskip

{
\textstyleWWPolicepardfaut{\textit{Quatrième étape
:}}\textstyleWWPolicepardfaut{
}\textstyleCodeInsr{indiceMédian}\textstyleWWPolicepardfaut{
}\textsf{←}\textstyleWWPolicepardfaut{ }\textstyleCodeInsr{(13 + 14)
}\textstyleCodeInsr{DIV}\textstyleCodeInsr{
2}\textstyleWWPolicepardfaut{, c’est-à-dire 13 où se trouve la valeur
cherchée ; la recherche est terminée.}}

{\sffamily\bfseries\upshape
Exemple de recherche infructueuse}

{
{Recherchons à présent la valeur
}{\textbf{8}}{. Les
premières étapes vont donner}}

\begin{flushleft}
\tablehead{}
\begin{supertabular}{m{0.507cm}m{0.509cm}m{0.509cm}m{0.509cm}m{0.507cm}m{0.509cm}m{0.509cm}m{0.509cm}m{0.509cm}m{0.507cm}m{0.509cm}m{0.509cm}m{0.509cm}m{0.509cm}m{0.51100004cm}m{0.51100004cm}m{0.509cm}m{0.51100004cm}m{0.51100004cm}m{0.527cm}}
\centering \itshape 1 &
\centering \itshape 2 &
\centering \itshape 3 &
\centering \itshape 4 &
\centering \itshape 5 &
\centering \itshape 6 &
\centering \itshape 7 &
\centering \itshape 8 &
\centering \itshape 9 &
\centering \itshape 10 &
\centering \itshape 11 &
\centering \itshape 12 &
\centering \itshape 13 &
\centering \itshape 14 &
\centering \itshape 15 &
\centering \itshape 16 &
\centering \itshape 17 &
\centering \itshape 18 &
\centering \itshape 19 &
\centering\arraybslash \itshape 20\\\hline
\multicolumn{1}{|m{0.507cm}|}{ 1} &
\multicolumn{1}{m{0.509cm}|}{ 3} &
\multicolumn{1}{m{0.509cm}|}{ 5} &
\multicolumn{1}{m{0.509cm}|}{ 7} &
\multicolumn{1}{m{0.507cm}|}{ 7} &
\multicolumn{1}{m{0.509cm}|}{ 9} &
\multicolumn{1}{m{0.509cm}|}{ 9} &
\multicolumn{1}{m{0.509cm}|}{ 10} &
\multicolumn{1}{m{0.509cm}|}{ 10} &
\multicolumn{1}{m{0.507cm}|}{ 15} &
\multicolumn{1}{m{0.509cm}|}{ 16} &
\multicolumn{1}{m{0.509cm}|}{ 20} &
\multicolumn{1}{m{0.509cm}|}{ 23} &
\multicolumn{1}{m{0.509cm}|}{ 23} &
\multicolumn{1}{m{0.51100004cm}|}{ 25} &
\multicolumn{1}{m{0.51100004cm}|}{ 28} &
\multicolumn{1}{m{0.509cm}|}{ 28} &
\multicolumn{1}{m{0.51100004cm}|}{ 28} &
\multicolumn{1}{m{0.51100004cm}|}{ 29} &
\multicolumn{1}{m{0.527cm}|}{ 29}\\\hline
\end{supertabular}
\end{flushleft}

\bigskip

\begin{flushleft}
\tablehead{}
\begin{supertabular}{m{0.507cm}m{0.509cm}m{0.509cm}m{0.509cm}m{0.507cm}m{0.509cm}m{0.509cm}m{0.509cm}m{0.509cm}m{0.507cm}m{0.509cm}m{0.509cm}m{0.509cm}m{0.509cm}m{0.51100004cm}m{0.51100004cm}m{0.509cm}m{0.51100004cm}m{0.51100004cm}m{0.527cm}}
\centering \itshape 1 &
\centering \itshape 2 &
\centering \itshape 3 &
\centering \itshape 4 &
\centering \itshape 5 &
\centering \itshape 6 &
\centering \itshape 7 &
\centering \itshape 8 &
\centering \itshape 9 &
\centering \itshape 10 &
\centering \itshape 11 &
\centering \itshape 12 &
\centering \itshape 13 &
\centering \itshape 14 &
\centering \itshape 15 &
\centering \itshape 16 &
\centering \itshape 17 &
\centering \itshape 18 &
\centering \itshape 19 &
\centering\arraybslash \itshape 20\\\hline
\multicolumn{1}{|m{0.507cm}|}{ 1} &
\multicolumn{1}{m{0.509cm}|}{ 3} &
\multicolumn{1}{m{0.509cm}|}{ 5} &
\multicolumn{1}{m{0.509cm}|}{ 7} &
\multicolumn{1}{m{0.507cm}|}{ 7} &
\multicolumn{1}{m{0.509cm}|}{ 9} &
\multicolumn{1}{m{0.509cm}|}{ 9} &
\multicolumn{1}{m{0.509cm}|}{ 10} &
\multicolumn{1}{m{0.509cm}|}{ 10} &
\multicolumn{1}{m{0.507cm}|}{ 15} &
\multicolumn{1}{m{0.509cm}|}{ 16} &
\multicolumn{1}{m{0.509cm}|}{ 20} &
\multicolumn{1}{m{0.509cm}|}{ 23} &
\multicolumn{1}{m{0.509cm}|}{ 23} &
\multicolumn{1}{m{0.51100004cm}|}{ 25} &
\multicolumn{1}{m{0.51100004cm}|}{ 28} &
\multicolumn{1}{m{0.509cm}|}{ 28} &
\multicolumn{1}{m{0.51100004cm}|}{ 28} &
\multicolumn{1}{m{0.51100004cm}|}{ 29} &
\multicolumn{1}{m{0.527cm}|}{ 29}\\\hline
\end{supertabular}
\end{flushleft}

\bigskip

\begin{flushleft}
\tablehead{}
\begin{supertabular}{m{0.507cm}m{0.509cm}m{0.509cm}m{0.509cm}m{0.507cm}m{0.509cm}m{0.509cm}m{0.509cm}m{0.509cm}m{0.507cm}m{0.509cm}m{0.509cm}m{0.509cm}m{0.509cm}m{0.51100004cm}m{0.51100004cm}m{0.509cm}m{0.51100004cm}m{0.51100004cm}m{0.527cm}}
\centering \itshape 1 &
\centering \itshape 2 &
\centering \itshape 3 &
\centering \itshape 4 &
\centering \itshape 5 &
\centering \itshape 6 &
\centering \itshape 7 &
\centering \itshape 8 &
\centering \itshape 9 &
\centering \itshape 10 &
\centering \itshape 11 &
\centering \itshape 12 &
\centering \itshape 13 &
\centering \itshape 14 &
\centering \itshape 15 &
\centering \itshape 16 &
\centering \itshape 17 &
\centering \itshape 18 &
\centering \itshape 19 &
\centering\arraybslash \itshape 20\\\hline
\multicolumn{1}{|m{0.507cm}|}{ 1} &
\multicolumn{1}{m{0.509cm}|}{ 3} &
\multicolumn{1}{m{0.509cm}|}{ 5} &
\multicolumn{1}{m{0.509cm}|}{ 7} &
\multicolumn{1}{m{0.507cm}|}{ 7} &
\multicolumn{1}{m{0.509cm}|}{ 9} &
\multicolumn{1}{m{0.509cm}|}{ 9} &
\multicolumn{1}{m{0.509cm}|}{ 10} &
\multicolumn{1}{m{0.509cm}|}{ 10} &
\multicolumn{1}{m{0.507cm}|}{ 15} &
\multicolumn{1}{m{0.509cm}|}{ 16} &
\multicolumn{1}{m{0.509cm}|}{ 20} &
\multicolumn{1}{m{0.509cm}|}{ 23} &
\multicolumn{1}{m{0.509cm}|}{ 23} &
\multicolumn{1}{m{0.51100004cm}|}{ 25} &
\multicolumn{1}{m{0.51100004cm}|}{ 28} &
\multicolumn{1}{m{0.509cm}|}{ 28} &
\multicolumn{1}{m{0.51100004cm}|}{ 28} &
\multicolumn{1}{m{0.51100004cm}|}{ 29} &
\multicolumn{1}{m{0.527cm}|}{ 29}\\\hline
\end{supertabular}
\end{flushleft}

\bigskip

\begin{flushleft}
\tablehead{}
\begin{supertabular}{m{0.507cm}m{0.509cm}m{0.509cm}m{0.509cm}m{0.507cm}m{0.509cm}m{0.509cm}m{0.509cm}m{0.509cm}m{0.507cm}m{0.509cm}m{0.509cm}m{0.509cm}m{0.509cm}m{0.51100004cm}m{0.51100004cm}m{0.509cm}m{0.51100004cm}m{0.51100004cm}m{0.527cm}}
\centering \itshape 1 &
\centering \itshape 2 &
\centering \itshape 3 &
\centering \itshape 4 &
\centering \itshape 5 &
\centering \itshape 6 &
\centering \itshape 7 &
\centering \itshape 8 &
\centering \itshape 9 &
\centering \itshape 10 &
\centering \itshape 11 &
\centering \itshape 12 &
\centering \itshape 13 &
\centering \itshape 14 &
\centering \itshape 15 &
\centering \itshape 16 &
\centering \itshape 17 &
\centering \itshape 18 &
\centering \itshape 19 &
\centering\arraybslash \itshape 20\\\hline
\multicolumn{1}{|m{0.507cm}|}{ 1} &
\multicolumn{1}{m{0.509cm}|}{ 3} &
\multicolumn{1}{m{0.509cm}|}{ 5} &
\multicolumn{1}{m{0.509cm}|}{ 7} &
\multicolumn{1}{m{0.507cm}|}{ 7} &
\multicolumn{1}{m{0.509cm}|}{ 9} &
\multicolumn{1}{m{0.509cm}|}{ 9} &
\multicolumn{1}{m{0.509cm}|}{ 10} &
\multicolumn{1}{m{0.509cm}|}{ 10} &
\multicolumn{1}{m{0.507cm}|}{ 15} &
\multicolumn{1}{m{0.509cm}|}{ 16} &
\multicolumn{1}{m{0.509cm}|}{ 20} &
\multicolumn{1}{m{0.509cm}|}{ 23} &
\multicolumn{1}{m{0.509cm}|}{ 23} &
\multicolumn{1}{m{0.51100004cm}|}{ 25} &
\multicolumn{1}{m{0.51100004cm}|}{ 28} &
\multicolumn{1}{m{0.509cm}|}{ 28} &
\multicolumn{1}{m{0.51100004cm}|}{ 28} &
\multicolumn{1}{m{0.51100004cm}|}{ 29} &
\multicolumn{1}{m{0.527cm}|}{ 29}\\\hline
\end{supertabular}
\end{flushleft}

\bigskip

{
Arrivé à ce stade, la zone de recherche s’est réduite à un seul élément.
Ce n’est pas celui qu’on recherche mais c’est à cet endroit qu’il se
serait trouvé ; c’est donc là qu’on pourra éventuellement
l'insérer. Le paramètre de sortie prend la valeur 6.}

\section{Introduction à la complexité}
{
L’algorithme de recherche dichotomique est beaucoup plus rapide que
l’algorithme de recherche linéaire. Mais qu’est-ce que cela veut dire
exactement ? Comment le mesure-t-on ? Quels critères utilise-t-on ?}

{
On pourrait se dire que pour comparer deux algorithmes, il suffit de les
traduire tous les deux dans un langage de programmation, de les
exécuter et de comparer leur temps d’exécution. Cela pose toutefois
quelques problèmes}

\liststyleListv
\begin{itemize}
\item {
Il faut que les programmes soient exécutés dans des environnements
strictement identiques ce qui n’est pas toujours facile}
\item {
Le test ne porte que sur un (voir quelques uns) jeu(x) de tests. Comment
en tirer un enseignement général ? Que se passerait-il avec des données
plus importantes ? Avec des données différentes ?}
\end{itemize}
{
En fait, un chiffre précis ne nous intéresse pas. Ce qui est intéressant
à connaitre, c’est l’évolution de la rapidité de l’algorithme (plus
précisément, le nombre d’opérations de base) en fonction de la
\textbf{taille} du problème à résoudre. Comment va-t-il se comporter
sur de «~gros~» problèmes ?}

{
Dans le cas de la recherche, la taille du problème est la taille du
tableau dans lequel on recherche un élément. On peut considérer que
l’opération de base est la comparaison avec un élément du tableau. La
question est alors la suivante : pour un tableau de taille N, à combien
de comparaisons faut-il procéder pour trouver l’élément (ou se rendre
compte de son absence) ?}

{\bfseries
Pour la recherche linéaire}

{
Cela dépend évidemment de la position de la valeur à trouver. Dans le
meilleur des cas c’est 1, dans le pire c’est N mais on peut dire, qu’en
moyenne, cela entraine «~N/2~» comparaisons (que ce soit pour la
trouver ou se rendre compte de son absence).}

{\bfseries
Pour la recherche dichotomique}

{
Ici, la zone de recherche est divisée par 2, à chaque étape. Imaginons
une liste de 64 éléments : après 6 étapes, on arrive à un élément
isolé. Pour une liste de taille N, on peut en déduire que le nombre de
comparaisons est au maximum l'exposant
qu'il faut donner à 2 pour obtenir N, soit
«~log\textsubscript{2 }N~».}

{\sffamily\bfseries\upshape
Comparaisons}

{
Voici un tableau comparatif du nombre de comparaison en fonction de la
taille N}

\begin{center}
\tablehead{}
\begin{supertabular}{|m{4.4300003cm}|m{0.507cm}|m{0.71900004cm}|m{0.93100005cm}|m{1.2479999cm}|m{1.4599999cm}|m{1.671cm}|}
\hline
\raggedleft \bfseries N &
\raggedleft \bfseries 10 &
\raggedleft \bfseries 100 &
\raggedleft \bfseries 1000 &
\raggedleft \bfseries 10.000 &
\raggedleft \bfseries 100.000 &
\raggedleft\arraybslash \bfseries 1
million\\\hline
\raggedleft \bfseries recherche linéaire &
\raggedleft  5 &
\raggedleft  50 &
\raggedleft  500 &
\raggedleft  5.000 &
\raggedleft  50.000 &
\raggedleft\arraybslash  500.000\\\hline
\raggedleft \bfseries recherche dichotomique &
\raggedleft  4 &
\raggedleft  7 &
\raggedleft  10 &
\raggedleft  14 &
\raggedleft  17 &
\raggedleft\arraybslash  20\\\hline
\end{supertabular}
\end{center}
{
On voit que c’est surtout pour des grands problèmes que la recherche
dichotomique montre toute son efficacité. On voit aussi que ce qui est
important pour mesurer la complexité c'est
l'ordre de grandeur du nombre de comparaisons. }

{
On dira que la recherche simple est un algorithme de complexité linéaire
(c'est-à-dire que le nombre
d'opérations est de l'ordre de N ou
proportionnel à N) ce qu’on note en langage plus mathématique O(N)
(prononcé «~grand O de N~»). Pour la recherche dichotomique, la
complexité est logarithmique, et on le note O( log\textsubscript{2 }N
).}

{
Comparons les complexités les plus courantes.}

\begin{center}
\tablehead{}
\begin{supertabular}{|m{3.167cm}|m{1.074cm}|m{1.43cm}|m{1.074cm}|m{1.4319999cm}|m{1.287cm}|m{1.425cm}|m{1.714cm}|}
\hline
\raggedleft \bfseries taille problème &
\raggedleft \bfseries 10 &
\raggedleft \bfseries 100 &
\raggedleft \bfseries 1000 &
\raggedleft \bfseries 10\textsuperscript{4} &
\raggedleft \bfseries 10\textsuperscript{5} &
\raggedleft \bfseries 10\textsuperscript{6} &
\raggedleft\arraybslash \bfseries
10\textsuperscript{9}\\\hline
\centering  O( 1 ) &
\raggedleft  1 &
\raggedleft  1 &
\raggedleft  1 &
\raggedleft  1 &
\raggedleft  1 &
\raggedleft  1 &
\raggedleft\arraybslash  1\\\hline
\centering  {O(
log}{\textsubscript{2
}}{N )} &
\raggedleft  4 &
\raggedleft  7 &
\raggedleft  10 &
\raggedleft  14 &
\raggedleft  17 &
\raggedleft  20 &
\raggedleft\arraybslash  30\\\hline
\centering  O( N ) &
\raggedleft  10 &
\raggedleft  100 &
\raggedleft  1000 &
\raggedleft  10.000 &
\raggedleft  10\textsuperscript{5} &
\raggedleft  10\textsuperscript{6} &
\raggedleft\arraybslash 
10\textsuperscript{9}\\\hline
\centering  O( N\textsuperscript{2} ) &
\raggedleft  100 &
\raggedleft  10.000 &
\raggedleft  10\textsuperscript{6} &
\raggedleft  10\textsuperscript{8} &
\raggedleft  10\textsuperscript{10} &
\raggedleft  10\textsuperscript{12} &
\raggedleft\arraybslash 
10\textsuperscript{18}\\\hline
\centering  O( N\textsuperscript{3} ) &
\raggedleft  1000 &
\raggedleft  10\textsuperscript{6} &
\raggedleft  10\textsuperscript{9} &
\raggedleft  10\textsuperscript{12} &
\raggedleft  10\textsuperscript{15} &
\raggedleft  10\textsuperscript{18} &
\raggedleft\arraybslash 
10\textsuperscript{27}\\\hline
\centering  O( 2\textsuperscript{N} ) &
\raggedleft  1024 &
\raggedleft  10\textsuperscript{30} &
\raggedleft  10\textsuperscript{301} &
\raggedleft  10\textsuperscript{3010} &
\raggedleft  10\textsuperscript{30102} &
\raggedleft  10\textsuperscript{301029} &
\raggedleft\arraybslash 
10\textsuperscript{301029995}\\\hline
\end{supertabular}
\end{center}
{
On voit ainsi que si on trouve un algorithme correct pour résoudre un
problème mais que cet algorithme est de complexité exponentielle alors
cet algorithme ne sert à rien en pratique. Par exemple un algorithme de
recherche de complexité exponentielle, exécuté sur une machine pouvant
effectuer un milliard de comparaisons par secondes, prendrait plus de
dix mille milliards d’années pour trouver une valeur dans un tableau de
100 éléments.}

{\sffamily\bfseries\upshape
Exercices}

\liststyleExercice
\begin{enumerate}
\item {\sffamily\bfseries
Calcul de complexités}
\end{enumerate}
\liststyleNumberingv
\begin{enumerate}
\item {
Quelle est la complexité d'un algorithme qui recherche
le maximum d'un tableau de N éléments ?}
\item {
Quelle est la complexité d'un algorithme qui remplace
par 0 toutes les occurrences du maximum d'un tableau
de N éléments ?}
\item {
Quelle est la complexité d'un algorithme qui vérifie si
un tableau contient deux éléments égaux ?}
\item {
Quelle est la complexité de l'algorithme de tri par
recherche des minima successifs ?}
\item {
Quelle est la complexité d'un algorithme qui vérifie si
un tableau est un palindrome ?}
\end{enumerate}
\liststyleExercice
\begin{enumerate}
\item {\sffamily\bfseries
Réflexion}
\end{enumerate}
{
L’algorithme de recherche dichotomique est-il toujours à préférer à
l’algorithme de recherche linéaire ?}

\section{Références}
\liststyleListv
\begin{itemize}
\item {
\url{http://interstices.info/jcms/c_6973/les-algorithmes-de-tri}

\textit{Ce site permet de visualiser les méthodes de tris en
fonctionnement. Il présente tous les tris vus en première plus quelques
autres comme le «~tri rapide~» («~quicksort~») que vous verrez en
deuxième année. }}
\end{itemize}
